\documentclass[10pt,a4paper]{article}
\usepackage[margin=20mm]{geometry}
\usepackage{fontspec}
\setmainfont{DejaVu Serif}
\usepackage{longtable,booktabs,array}
\usepackage[hidelinks]{hyperref}
\setlength{\parindent}{0pt}
\setlength{\parskip}{6pt}
\begin{document}
\begin{center}{\Large\bfseries Separating Numeric Support from Reading Competence}\par
WikiGraphRAG v17: preparation and failure-analysis checkpoint\par
4 October 2026\end{center}
\textit{Working addendum. Generic TeX source; ACL template and compilation are not verified.}
\section*{Abstract}
Exact source numbers do not establish correct report, period, entity or population binding. This development checkpoint separates conventional numeric support, qualified source views and reader competence. Pinned processor inventories cover 23,651 numeric occurrences and agree on all 23,519 commonly normalized values, while nine clean source outcomes, one incomplete dependency set and two telemetry failures remain distinct. All 101 original evidence blocks have qualified inspection views across preserved original and recovery attempts. Identity and whole-bundle admission remain incomplete. Two separately recorded native attempts yield no prompt tokenization: the first is refused before launch and the second stops during loading. These results support preparation and failure analysis, not natural question-answering gains, a stronger reader or a new retrieval algorithm.
\section*{1. A preparation checkpoint, not a reader result}
The research target is automatic question-to-evidence binding: selecting every quantity and its required source, period, unit, entity and population from an acquired pack. Exact numeric support is necessary but does not determine whether the selected binding answers a separately written request. The present work establishes conventional baselines and diagnoses execution and source-view failures before a new retrieval policy is evaluated.

The cohort remains the known development filings. Source-informed table selection and reference-informed complete evidence will make the planned reader experiment an oracle-evidence diagnostic, not retrieval success, human gold or held-out generalization. The earlier preparation addendum and every failed attempt remain unchanged.

\begin{longtable}{@{}p{0.2018\textwidth}p{0.3652\textwidth}p{0.3729\textwidth}@{}}\toprule
\textbf{Stage} & \textbf{Observed outcome} & \textbf{Interpretation} \\ \midrule\endhead
Pinned processor & Inventories saved; source statuses retained & Conventional numeric support baseline \\
Source views & Original and recovery outcomes preserved & Qualified views available; evidence closure separate \\
Identity and bundles & Candidate views and staging built & Complete identity cards and evidence remain unadmitted \\
Native reader & Prelaunch refusal; separate loading stop & No tokenizer or reading-quality measurement \\
Natural study & No evaluated reader outputs & Prospective workflow only \\
\bottomrule\end{longtable}
{\small These stages have different units and are not an accuracy leaderboard.}
This is a research preparation and failure-analysis checkpoint. It is not an ACL-ready empirical contribution; admission, natural inference and the matched reader comparison have not been completed.

\clearpage
\section*{2. The conventional processor baseline}
The new execution uses pinned Arelle with captured dependencies, explicit physical DOM traversal and the reviewed process-group monitor. Its predecessor attempts are retained: the original total-memory gate launched no worker, and the first cache-aware attempt exposed traversal and process-identity defects. The published intervening draft was not executed.

The incomplete source retains the single timed-out SIC dependency. The two telemetry failures occur when live-process RSS becomes unreadable near exit; they are not observed RSS-cap exceedances. Both save complete inventories, exit successfully and have no model or pin errors, but the frozen watchdog statuses stay failures.

\begin{longtable}{@{}p{0.3345\textwidth}p{0.1057\textwidth}p{0.2115\textwidth}p{0.2883\textwidth}@{}}\toprule
\textbf{Preserved source status} & \textbf{Sources} & \textbf{Occurrences} & \textbf{Exact common agreements} \\ \midrule\endhead
Clean under enabled checks & 9 & 16064 & 15984 \\
Incomplete dependency set & 1 & 2905 & 2900 \\
RSS telemetry failure & 2 & 4682 & 4635 \\
All saved inventories & 12 & 23651 & 23519 \\
\bottomrule\end{longtable}
{\small Common agreements compare the overlapping non-nil normalized population. Saved output does not turn a retained technical failure into successful validation.}
\begin{longtable}{@{}p{0.6517\textwidth}p{0.2883\textwidth}@{}}\toprule
\textbf{Inventory property} & \textbf{Observed occurrences} \\ \midrule\endhead
Previously unsupported SEC number-word forms normalized & 116 \\
Both processors report nil & 16 \\
Effective English labels available & 23651 \\
Unresolved numeric concepts & 0 \\
Recorded discrepancy locators & 0 \\
\bottomrule\end{longtable}
{\small Normalization and label lookup are observed capabilities of the pinned conventional processor; they are not new algorithms.}
The requested core inline and taxonomy checks exclude SEC filing rules, calculation validation, UTR and formulas. Agreement on the same acquired filing bytes does not establish SEC-original byte identity, financial truth, complete conformance or question-answering quality.

\clearpage
\section*{3. Source views and identity require separate admission}
Office conversion initially produces source views for most requested blocks. Subsequent all-page review checks glyphs, clipping, numeric-column mapping and original CSS. These are disjoint model-role source subsets, not duplicate independent annotation or an inter-annotator-agreement experiment. The original conversion and visual-exclusion outcomes remain fixed.

Separate recovery first handles retained telemetry stops, then a source-specific contrast defect. The contrast amendment changes only a narrow CSS alpha notation under a byte-level patch ledger; original source text and locators remain unchanged. Authored controls support the tested text/background and merged-column behavior, while low-alpha border fidelity still fails. The views remain qualified office-rendered inspection aids.

\begin{longtable}{@{}p{0.6517\textwidth}p{0.2883\textwidth}@{}}\toprule
\textbf{Original request or output} & \textbf{Blocks} \\ \midrule\endhead
Selected tables & 21 \\
Unique existing neighbors & 80 \\
Total requested blocks & 101 \\
Converted and rasterized & 98 \\
Retained telemetry stops & 3 \\
\bottomrule\end{longtable}
{\small The original conversion denominator is fixed. Recovery and extra identity dependencies require separately bound records; they never replace these outcomes.}
\begin{longtable}{@{}p{0.6517\textwidth}p{0.2883\textwidth}@{}}\toprule
\textbf{Qualified-view availability stage} & \textbf{Blocks or pages} \\ \midrule\endhead
Original-attempt blocks admitted & 90 \\
Separate telemetry-recovery blocks admitted & 3 \\
Separate contrast-recovery blocks admitted & 8 \\
Original-population blocks available across attempts & 101 \\
Actual raster pages inspected across these attempts & 109 \\
\bottomrule\end{longtable}
{\small The block population stays fixed; repeated rendering creates additional inspected pages. Post-inspection repair outcomes never replace earlier exclusions.}
The fixed original population now has qualified views for every selected table and existing neighbor across separate attempts. Source text and exact anchors are retained. This does not establish browser equivalence, complete supporting evidence or correct semantic attachment.

\clearpage
\section*{4. Two retained native feasibility failures}
The reviewed recipe permits one server load, health/identity inspection and authored-prompt tokenization per frozen attempt, with no generation or natural prompt. The first attempt is refused before asset verification or process launch because its pressure proxy plus owned-process-group reserve exceeds the threshold. A separately approved follow-up, after exact diagnostic-file deduplication, keeps the same model, runtime, profile, caps and authored prompts.

The follow-up verifies the pinned assets and starts partial loading. The watchdog then kills the owned group when shared cgroup memory.current reaches the unchanged hard-limit value. All health polls fail; final cleanup confirms that the owned group is absent. Neither attempt establishes full model compatibility, cache allocation, tokenizer fit or reading competence.

\begin{longtable}{@{}p{0.421\textwidth}p{0.2595\textwidth}p{0.2595\textwidth}@{}}\toprule
\textbf{Attempt measure} & \textbf{v17: prelaunch refusal} & \textbf{v17.1: loading stop} \\ \midrule\endhead
Planned authored tokenizations & 8 & 8 \\
Native processes started & 0 & 1 \\
Prompts tokenized & 0 & 0 \\
Completion calls & 0 & 0 \\
\bottomrule\end{longtable}
{\small These are separate frozen attempts. Neither passes health/identity or produces a native tokenizer or performance measurement; every planned prompt remains unattempted.}
\begin{longtable}{@{}p{0.6324\textwidth}p{0.3076\textwidth}@{}}\toprule
\textbf{Follow-up loading observation} & \textbf{Bytes or events} \\ \midrule\endhead
Peak shared cgroup total (bytes) & 8589934592 \\
Peak pressure proxy (bytes) & 5057638400 \\
Peak owned-group RSS (bytes) & 2326736896 \\
Observed OOM events (increment) & 0 \\
Observed OOM-kill events (increment) & 0 \\
\bottomrule\end{longtable}
{\small The application equality-stop terminates loading at the shared total boundary. This is not an observed kernel OOM kill or a measured standalone model-memory requirement.}
The conservative pressure proxy and owned-group RSS remain below their separate caps at the follow-up stop. Kernel reclaim can operate at the hard limit (Linux Kernel documentation); equality alone is not proof of OOM. File-cache pressure is a plausible explanation, not an established cause or evidence that the full model fits. A guard-only proposal would retain the hard limit, proxy/reserve, RSS, address-space, time and OOM-event checks while removing the application equality stop. It remains unexecuted, with shared-cgroup risk explicit.

\clearpage
\section*{5. Identity and whole-bundle preparation}
A provenance audit found that legacy identity cards copied DEI values without hiding flags. Replacement cards require visible issuer, report-type and report-date support; calendar report-end years cannot substitute for fiscal-year labels. One NVDA report-type view retains its telemetry failure. The separate SLB cover views support literal visible spelling only: a proprietary hidden-fact style reference and legal-name equivalence remain unverified.

The source-conditioned bundle inventory adds complete heading, definition and continuation blocks while retaining all existing neighbors and acquired numeric occurrences. Copied concept, entity, period and unit aspects are declared machine metadata, not visible semantic certification. Every population remains null, and the additional dependency views still require rendering and review.

\begin{longtable}{@{}p{0.6517\textwidth}p{0.2883\textwidth}@{}}\toprule
\textbf{Identity preparation measure} & \textbf{Count} \\ \midrule\endhead
Required issuer/report/date fields & 36 \\
Fields with selector candidates & 34 \\
Identity candidate views requested & 34 \\
Identity candidate views admitted & 33 \\
Separate SLB literal-cover views admitted & 4 \\
Complete source identity cards constructed & 0 \\
\bottomrule\end{longtable}
{\small The identity and SLB populations are separate from the original evidence blocks. Qualified views do not imply legal-name or hidden-metadata equivalence.}
\begin{longtable}{@{}p{0.6517\textwidth}p{0.2883\textwidth}@{}}\toprule
\textbf{Question-free preparation measure} & \textbf{Count} \\ \midrule\endhead
Fixed table bundles materialized & 21 \\
Additional dependency blocks identified & 40 \\
Physically separated source-only projections & 21 \\
Complete evidence bundles admitted & 0 \\
Natural questions authored & 0 \\
Natural references authored & 0 \\
\bottomrule\end{longtable}
{\small Dependency suggestions and schema checks precede semantic admission. The new blocks are a separate population and are not added to the original source-view denominator.}
Ten candidates pass interface-shape validation; eleven retain block\_text failures from mandatory blank neighbors. No neighbor is removed or schema changed. Physically separated source-only exports prevent handing authors the combined private file with typed records. Syntax and export checks admit no author, reference or complete evidence pack.

\clearpage
\section*{6. The decision before a method claim}
The planned lexical, model-proposal and direct-reader arms receive identical acquired evidence, including canonical numeric values. They therefore test machine-assisted selection and scope, not raw number decoding. Question meaning remains externally graded. Source-only authors and two fresh model-role reviewers have separate access; their annotations are model-generated development data, not human gold.

Control admission must remove every duplicate disclosure of the chosen relation while preserving quantity lexemes and values. A token serving both roles makes the control unavailable. Semantic graders receive the same final action, quantity, expanded scope and citation projection across arms; proposal/finalizer diagnostics remain separate. Fixed proposer-then-direct order and cold-server confounding are disclosed, with no latency-superiority claim. Realized beam sizes, invalid outputs and all unattempted cells remain in the denominator.

\begin{longtable}{@{}p{0.6517\textwidth}p{0.2883\textwidth}@{}}\toprule
\textbf{Stage A design quantity} & \textbf{Planned count} \\ \midrule\endhead
Separately authored base questions & 12 \\
Missing-evidence conditions & 6 \\
Total acquired-evidence conditions & 18 \\
System outputs across the three methods & 54 \\
Native completion calls across both model arms & 36 \\
\bottomrule\end{longtable}
{\small Design counts only: these are not admitted or executed results. The full planned denominator remains visible if later admission or execution fails.}
A future retrieval contribution must first show residual binding failures after competent complete-evidence reading. Equivalence classes and weighted cross-class edge cutting already appear in EC2 (Golovin et al., 2010); ECED treats correlated noisy tests (Chen et al., 2017). Budget-aware acquisition and adaptive stopping also need cost-aware RAG baselines (Wu et al., 2026). These precedents do not automatically supply guarantees for learned utility, incomplete beams or conjunctive evidence.

Reproducibility: public result digests and numerical table bindings accompany the canonical manuscript. Source text and private registries stay external under content hashes. The PDF uses ReportLab; the editable generic TeX is uncompiled and ACL template compliance is unverified. Prior attempts and the earlier addendum remain unchanged.

\section*{References}
Daniel Golovin, Andreas Krause and Debajyoti Ray. 2010. Near-Optimal Bayesian Active Learning with Noisy Observations. Advances in Neural Information Processing Systems 23. \url{https://proceedings.neurips.cc/paper/2010/hash/1e6e0a04d20f50967c64dac2d639a577-Abstract.html}

Yuxin Chen, Hamed Hassani and Andreas Krause. 2017. Near-optimal Bayesian Active Learning with Correlated and Noisy Tests. AISTATS, PMLR 54:223-231. \url{https://proceedings.mlr.press/v54/chen17b.html}

Mingyan Wu, Han Yang, Omer Ben-Porat and Yftah Ziser. 2026. When Knowledge Is Not Free: Cost-Aware Evidence Selection in Retrieval-Augmented Generation. arXiv:2606.02245v1. \url{https://arxiv.org/abs/2606.02245v1}

Linux Kernel documentation. Control Group v2: memory.max, memory.events and memory.stat. Accessed 4 October 2026. \url{https://docs.kernel.org/admin-guide/cgroup-v2.html}

\end{document}
